% Standalone source-record documentation; no research attempts.
\documentclass[11pt]{article}
\usepackage[margin=25mm]{geometry}
\usepackage{fontspec,unicode-math}
\setmainfont{Latin Modern Roman}
\setmathfont{Latin Modern Math}
\usepackage{amsmath,amssymb,mathtools,xurl,hyperref}
\hypersetup{colorlinks=true,urlcolor=blue,linkcolor=blue}
\setcounter{secnumdepth}{0}
\setlength{\parindent}{0pt}
\setlength{\parskip}{5pt}
\setlength{\emergencystretch}{3em}
\newcommand{\sref}[2]{\hyperlink{source:#1:#2}{[S#2]}}
\newcommand{\cataloguescope}[1]{\paragraph{Source record.}#1}
\newcommand{\floor}[1]{\left\lfloor#1\right\rfloor}
\begin{document}
% UnsolvedMath ID 76; source catalogue code DARPA-006
\setcounter{section}{75}
\section{Computational Duality}\label{76}
{\small\sffamily UnsolvedMath ID 76 · Computer Science}
\subsection{Definitions and mathematical statement}

\paragraph{Shared notation.}Standard number systems are $\mathbb N,\mathbb Z,\mathbb Q,\mathbb R,\mathbb C$. Specialized notation and index ranges follow the cited source and the definitions below.


A dual formulation associates another mathematical description with a given problem. For example, an optimization primal and dual compare minimization and maximization values; weak duality supplies an ordering, while strong duality is equality under appropriate hypotheses. A geometric representation uses structures such as convex sets, metric spaces or algebraic varieties. A computational advantage requires a specified task, input representation, correctness condition and resource comparison.
\paragraph{Record scope and qualifications.}
The catalogue title is retained as metadata. Named certificates, contracts, bounds or frameworks in that title are not asserted as results of the original source extract.
This is DARPA Mathematical Challenge Six. The source does not select a duality, input class, computational resource or baseline algorithm. Optimization duality is an explanatory example, rather than the only intended interpretation. The catalogue-linked mirror returned HTTP 404 during this preparation; its full original solicitation was not reverified. The supplied catalogue extract and historical attribution are documented with that access limitation.
\paragraph{Statement recorded in the supplied source.}
Use mathematical duality and geometry as foundations for developing novel computational algorithms. \sref{76}{2}
\paragraph{Source status.}
The supplied catalogue labels this source record open; that metadata was checked against the supplied snapshot on 2026-10-11. The label is preserved for provenance and is not an independent verification of present open status or mathematical correctness.
\subsection{Short English statement}
Investigate whether duality and geometric descriptions can lead to principled new algorithms.
\subsection{Sources}
\begin{itemize}
\item[\hypertarget{source:76:1}{S1}] Ulam AI and the UnsolvedMath Contributors, UnsolvedMath: A Curated Collection of Open Mathematics Problems, record 76 (DARPA-006). Catalogue curation and metadata: CC BY 4.0. The source extract is retained; definitions, English summaries and qualifications are editorial additions.
\url{https://huggingface.co/datasets/ulamai/UnsolvedMath}.
\item[\hypertarget{source:76:2}{S2}] DARPA DSO, Broad Agency Announcement 07-68, Mathematical Challenges (2007), challenge identified above; the catalogue-linked University of Tennessee mirror.
\url{https://web.math.utk.edu/~vasili/refs/darpa07.MathChallenges.html}.
\end{itemize}
\paragraph{Preparation and review.}
The original source and UnsolvedMath contributors retain attribution for the source extract. Codex (GPT-6; exact serving revision not exposed) assisted with definitions, English summaries, scope notes and formatting. No new human mathematical authorship or independently verified solution is asserted. The supplied catalogue credits: DARPA. Permanent source ID 76 and code DARPA-006 checked against the complete supplied registry and current repository. Distinct upstream records, including repeated questions and broad or damaged extracts, are retained for complete coverage of the source CS category. A repeated record is not a new independent conjecture. Complete mathematical-equivalence and current-literature review has not been performed. Preparation and final review follow contribution prompts C and D, with source gaps disclosed. Structural and rendering checks do not establish a new theorem or open status.
\label{end:76}
\end{document}
